\documentclass{article}

% --- PAQUETS DE BASE ---
\usepackage[utf8]{inputenc}
\usepackage[T1]{fontenc}
\usepackage{graphicx}       % Pour insérer des images
\usepackage[french]{babel}  % Pour les conventions typographiques françaises

% --- PAQUETS POUR LE CODE ET LES LÉGENDES ---
\usepackage{listings}       % Pour insérer des blocs de code
\usepackage{caption}        % Pour personnaliser les légendes
\usepackage{xcolor}         % Pour utiliser des couleurs personnalisées

% --- CONFIGURATION ---

% 1. Configuration du style des légendes (votre demande initiale, corrigée)
\captionsetup{
    justification=centering,  % Centrer la légende
    singlelinecheck=false,    % Appliquer la justification même si la légende tient sur une ligne
    font=small,               % Police plus petite pour la légende
    labelfont=bf              % "Figure 1", "Listing 1" en gras
}

% 2. Configuration de l'apparence des blocs de code (avec `listings`)
\lstdefinestyle{mystyle}{
    backgroundcolor=\color{black!5},   % Fond gris très clair
    commentstyle=\color{green!50!black},
    keywordstyle=\color{blue},
    numberstyle=\tiny\color{gray},
    stringstyle=\color{purple},
    basicstyle=\ttfamily\footnotesize, % Police à chasse fixe (monospace) et petite
    breakatwhitespace=false,         
    breaklines=true,                 
    captionpos=b,                    % Position de la légende en bas
    keepspaces=true,                 
    numbers=left,                    % Numéros de ligne à gauche
    numbersep=5pt,                   
    showspaces=false,                
    showstringspaces=false,
    showtabs=false,                  
    tabsize=2,
    frame=single,                    % Encadrer le code
    rulecolor=\color{black}          % Couleur du cadre
}

\lstset{style=mystyle} % Appliquer ce style par défaut à tous les listings

% --- DOCUMENT ---
\begin{document}

\section{Code Java}


\begin{figure}[h!]
    \centering
    \includegraphics[width=0.8\textwidth]{figures/code_java/default/default/fig_1.png}
    \caption{Caption here}
    \label{fig:code_java_default_1}
\end{figure}


\section{Autre codes}

Ceci est un exemple de code Java affiché correctement dans un document LaTeX.

% Utilisation de l'environnement lstlisting
\begin{lstlisting}[language=Java, caption={Exemple "Hello, World!" en Java.}, label={lst:hello_java}]
/* Code en Java */
public class Main {
  public static void main(String[] args) {
    System.out.println("Hello, World!");
  }
}
/* Fin du code */
\end{lstlisting}

Nous pouvons faire référence à ce code plus loin dans le texte en utilisant son étiquette, comme ceci : voir le listing \ref{lst:hello_java}.

\end{document}








